\section{Background and Related Work}\label{sec:background}

Kinodynamic planning involves navigating the complex interplay between the robot's \textsl{state space}---comprising position, velocity, and higher-order derivatives---and its \textsl{control space}, defined by the set of feasible actions for a given task.
The existing literature mainly bifurcates this analysis into two approaches: \textsl{sequential} and \textsl{interleaved}. 
Sequential methods, such as TrajOpt \cite{schulman2013finding}, CHOMP \cite{ratliff2009chomp}, STOMP \cite{kalakrishnan2011stomp}, KOMO \cite{toussaint2014newton}, and ITOMP \cite{park2012itomp}, 
construct a geometric path in the robot's state space first and then assign valid controls to each waypoint. %Exemplified by the likes of
While these methods are advantageous for their minimal design complexity and capability to handle non-linear constraints, they have a number of limitations.
They are highly susceptible to the quality of the initial seed trajectories, often confining trajectories to a narrow set homotopic classes \cite{schulman2013finding,ichnowski2020deep}.
Moreover, they generally lack theoretical guarantees regarding the completeness or optimality of solutions \cite{kuntz2019fast}.
Within this line of work, there exist hybrid methods that attempt to reconcile these limitations by incorporating sampling-based planning, which in turn brings about theoretical properties such as probabilistic completeness and asymptotic optimality \cite{kuntz2019fast} and permits the discovery of multiple solution classes \cite{kamat2022bitkomo}.
Time parameterization techniques can also be applied in place of optimization to assign timesteps to each waypoint in the geometric path \cite{pham2014general,Berscheid-RSS-21,macfarlane}.

In contrast to sequential approaches, interleaved methods cast motion planning as a search problem, alternating between state space sampling and control propagation to find a valid trajectory.
These methods often employ steering functions, or inverse models, to delineate a set of controls connecting two arbitrary points in the state space. 
For instance, DIMT-RRT utilizes a quadratic steering function to account for joint acceleration limits and non-zero initial and goal velocities \cite{kunz2014probabilistically}.
AVP computes the range of reachable velocities at the end of a path, given reachable velocities at the start of it \cite{pham2013kinodynamic,johnson2012optimal}.
LQR-RRT* uses a quadratic cost function and linearized dynamics in conjunction with a linear-quadratic regulator (LQR) to connect consecutive states in an RRT* tree \cite{perez2012lqr,caldwell2018fast}.
Notably, the optimization techniques discussed in sequential methods can also be used as steering functions in interleaved approaches.
In all, these techniques enable precise goal state attainment, allow to perform backward searches, and facilitate exact tree connections in bidirectional searches, thus expediting the planning process \cite{lavalle2006planning}.
However, they come with a set of challenges. Steering functions may not exist or could be difficult to design for a given dynamical system. Furthermore, exact state connections may not be reliable in real-world settings due to uncertainties and the use of optimization as steering may cause planning to get stuck in local minima \cite{littlefield2020efficient}.

The second variety of interleaved methods involves navigating the state and control spaces using forward control propagation. Using a random state and action, these methods determine the resultant state, constructing a tree of dynamically feasible edges \cite{pasricha2022pokerrt}.
SyCLoP and KPIECE leverage state space discretizations and guide exploration in low-coverage areas \cite{plaku2010motion,sucan2011sampling}.
SST optimizes path quality via pruning and heuristic biasing \cite{li2016asymptotically}.
GBRRT and GABRRT combine multiple propagation methods and backward tree costs as heuristic values for bidirectional kinodynamic search \cite{nayak2022bidirectional}. Other techniques involve pre-mapping the space of valid motions and conducting a search over this state and control space discretization \cite{shome2021asymptotically,bordalba2018randomized}.
Existing approaches that exploit lazy strategies for collision checking involve learning-based methods \cite{yu2021reducing} and evaluating paths on an as-needed basis \cite{kavraki2000path,mandalika2019generalized,hauser2015lazy}. We take an empirical approach by computing a collision probability for each kinodynamic action in our state and control space discretization.

Overall, interleaved approaches are easy to design since they leverage general-purpose physics simulations, explore both state and control spaces uniformly, and easily adapt to environmental interactions and multi-object systems. However, they may face slow convergence owing to costly simulation and collision checking methods. Next, we outline how our contribution mitigates these issues, paving the way for fast computation of asymptotically optimal solutions.

